\section{What would actually imply quasi-polynomial recognition?}
\label{sec:complexity}

The new results remove a local combinatorial obstruction to efficient
normalization.  They do not bound the number, size, or required geometry
of the states visited by a complete knot algorithm.  This section makes
the distinction quantitative, so that future work has a falsifiable
complexity target rather than an analogy with a fast local benchmark.

\subsection{A conditional composition theorem}

\begin{proposition}[Explicit accounting for a complete implementation]
\label{prop:composition}
Suppose a correct total unknot-recognition implementation on diagrams of
size $n$ makes at most $Q(n)$ calls to its geometric state-processing
operator.  Suppose all state encodings have size at most $N_*(n)$, every
exceptional interface has size at most $U(n)$, and all remaining work per
state, including certificate generation, geometric transport, and replay,
is bounded by a fixed polynomial in $N_*(n)+1$.  Allow at most
$(N_*(n)+1)^{O(1)}$ selector invocations per state, each on an input of
encoded size $(N_*(n)+1)^{O(1)}$.  Assume also that all work outside these
state calls, including preprocessing and finalization, is polynomial in
$n+N_*(n)+1$.  If it uses the exact finite selector of
Theorem~\ref{thm:exceptional-selector}, then its total time is bounded by
\begin{equation}
 T(n)\le (Q(n)+1)\,(n+N_*(n)+1)^{O(1)}
       2^{O(U(n)\log(U(n)+2))}.
 \label{eq:composition}
\end{equation}
In particular, the three bounds
\begin{align*}
 \log(Q(n)+2)&=O((\log(n+2))^2),\\
 \log(N_*(n)+2)&=O((\log(n+2))^2),\\
 U(n)\log(U(n)+2)&=O((\log(n+2))^2)
\end{align*}
suffice for $T(n)=n^{O(\log n)}$ for $n\ge2$.
\end{proposition}
\begin{proof}
Each finite selection costs at most a polynomial in its encoding times
the Bell-number factor.  The input-size hypothesis and the bound on the
number of invocations absorb all selector calls into the same fixed
polynomial factor per state.  Sum this and the remaining per-state costs
over at most $Q(n)$ calls, then absorb the outside work using $Q(n)+1$
and the polynomial in $n+N_*(n)+1$.  Taking logarithms gives terms bounded
by $\log(Q(n)+2)$, $O(\log(n+N_*(n)+1))$, and
$O(U(n)\log(U(n)+2))$.  Each is $O((\log n)^2)$ under the stated
hypotheses.  Exponentiating gives the conclusion.
\end{proof}

This is complexity accounting, not a new total-recognition theorem.
The first sentence assumes correctness and termination of a complete
implementation.  Our existing partial census procedure does not satisfy
that premise merely because every emitted certificate is correct.

A more convenient sufficient state-size condition is $N_*(n)=n^{O(1)}$.
Under this stronger condition, $U=O(\log n)$ gives local cost
\[
 2^{O(\log n\log\log n)}=n^{O(\log\log n)}.
\]
The weaker interface bound
$U=O((\log n)^2/\log\log n)$ also suffices for the target local
$n^{O(\log n)}$ estimate, away from small $n$.  Neither interface bound
is proved for arbitrary native inputs or along the required hierarchies.
The absence of a dependence exponential in the number of top handles
does not itself control $U$.

\subsection{Hierarchy depth, loop counts, and the displayed exponent}

The loop estimate $L(g+1)^L$ from Proposition~9.1 of
\cite{Lackenby2026} uses the maximum hierarchy length $L$ and the maximum
\emph{pattern complexity} $g$ of its decomposing surfaces.  This complexity
includes boundary-pattern intersections; a genus bound alone does not
control it.  The estimate immediately shows why quantitative depth
matters.  If $g\le n^a$ and $L\le b\log n$, then
\[
 \log\bigl(L(g+1)^L\bigr)
 =\log L+L\log(g+1)=O((\log n)^2).
\]
If instead $L=O((\log n)^2)$ with the same pattern-complexity bound, this calculation
gives $O((\log n)^3)$ in the exponential, namely
$n^{O((\log n)^2)}$.  A bound merely linear in a handle-complexity
parameter that can itself be polynomial in $n$ does not give either
logarithmic-depth conclusion.

The purpose of comparing the two 2021 slide estimates is therefore
specific: an implementation and proof should state which exponent they
establish.  Replacing a polynomial local factor by a faster polynomial
does not repair an uncontrolled hierarchy depth.  Conversely, a
quasi-polynomial number of states can still tolerate another
quasi-polynomial local factor of the same logarithmic exponent, as
Proposition~\ref{prop:composition} shows.

\subsection{Encoding is part of the theorem}

Let $t$ count the source tetrahedra, and let $b$ be the total bit length
of the supplied normal coordinates.  A local source-to-certificate bound
should be expressed in $t+b$ and in the encoded marked interfaces it
creates.  It should not hide a scan through $W$ parallel sheets behind a
phrase such as ``construct the cut.''  The distinction is familiar from
the orbit-counting machinery of Agol--Hass--Thurston
\cite{AgolHassThurston}; its algorithmic value is that large multiplicities
can be processed through finite interval data.

The deferred theorem never requires expanding those sheets: it operates
on a supplied finite incidence matrix.  But this does not prove that
every matrix needed later in a hierarchy will have small encoded size.
Annular paths, port refinements, boundary patterns, and repeated chart
substitutions can dominate the size even when scalar multiplicities
remain in binary.  Every successor operator therefore needs an explicit
bound on both the number of records and the bits in each record.

Output size creates an independent barrier.  If a cut has exponentially
many identical ball components, writing one certificate per component
cannot take polynomial time in a binary multiplicity.  The existing
census representation groups such components by genus and multiplicity.
A full hierarchy certificate needs an equally explicit policy for
sharing repeated local structures and for authenticating that sharing.

\subsection{Outstanding geometric and algorithmic obligations}

The known local hierarchy theorems cited in Section~\ref{sec:scope} are
valuable starting points.  The following are the stronger or
implementation-specific obligations still absent from this package.

\begin{enumerate}
\item \emph{Controlled progress.}  Produce the decomposing multisurfaces
or other accelerated progress objects needed for a short global
construction, and prove their depth and pattern-complexity bounds,
including boundary-pattern intersections, in the actual source encoding.
Finding an arbitrary nonseparating surface is not the open
subproblem here.
\item \emph{Source-authenticated successor transport.}  Implement the
applicable compressed cutting and parallelity constructions with explicit
marked boundary maps, preserving the hypotheses needed for the next
state.  A census is less information than a successor handle structure.
\item \emph{Repair and restart accounting.}  Bound all branches caused
by unsuccessful local candidates, changes of charts, boundary-pattern
repairs, or restarts.  Charging only successful cuts omits real work.
\item \emph{A complete terminal procedure.}  The raw certificate class
is not complete for all handlebodies in arbitrary charts.  A total
algorithm must explain what happens after a universal raw obstruction,
and bound the cost of its alternative.  It must also decide the relevant
boundary-pattern condition, not just whether the underlying pieces are
balls.
\item \emph{Exceptional-interface control.}  Bound $u$, replace it by a
more stable parameter, or find a polynomial selector for the actual
geometric class.  The abstract triangle witness prevents one particularly
simple invariance argument but does not settle that geometric question.
\end{enumerate}

These obligations interact.  A move that reduces $u$ while making marked
paths exponentially long may lose overall.  A normalization that reduces
the number of records but changes the admissible pattern may invalidate
the next theorem.  The intended measure must charge both geometry and
encoding, with a source verifier checking each claimed transition.

\subsection{Algebraic simplification and geometric transport}

Whitehead minimization has polynomial algorithms in the length of
explicitly written free-group words \cite{RoigVenturaWeil}.  Such a bound
is not automatically polynomial in the size of a compressed grammar for
those words.  More fundamentally, an algebraic automorphism does not by
itself give the marked meridian and port data expected by the native
geometric kernel.  The recent thickenable-presentation work
\cite{LackenbyStable2026} provides relevant structural context, but no
bounded marked-slide implementation is supplied here.

Section~\ref{sec:terminal-free-group} gives a legitimate alternative
use: under authenticated irreducibility hypotheses, a verified free
fundamental group proves an abstract terminal handlebody.  It can bypass
the need to interpret each algebraic transformation as a slide when only
that terminal conclusion is required.  It still needs a bound on the
source presentation, the certificate, and its discovery.  It also cannot
be substituted for marked transport when the algorithm must continue.
